% Generated draft: requires mathematical and source review.
\documentclass[11pt]{article}
\usepackage[margin=25mm]{geometry}
\usepackage{fontspec}
\setmainfont{Latin Modern Roman}
\usepackage{amsmath,amssymb,mathtools,unicode-math}
\setmathfont{Latin Modern Math}
\usepackage{enumitem,needspace,xurl,hyperref}
\setlength{\parindent}{0pt}
\setlength{\parskip}{5pt}
\setlength{\emergencystretch}{3em}
\newcommand{\N}{\mathbb{N}}
\newcommand{\Z}{\mathbb{Z}}
\newcommand{\Q}{\mathbb{Q}}
\newcommand{\R}{\mathbb{R}}
\newcommand{\C}{\mathbb{C}}
\newcommand{\F}{\mathbb{F}}
\newcommand{\A}{\mathbb{A}}
\newcommand{\PP}{\mathbb{P}}
\newcommand{\E}{\mathbb{E}}
\newcommand{\eps}{\varepsilon}
\newcommand{\dd}{\,\mathrm{d}}
\begin{document}
\section*{UnsolvedMath ID: 2200019}
\textbf{Source problem code:} catalogue entry
\subsection{Definitions and mathematical statement}
Let $p\in\R[x]$ have even degree, and let $p',p''$ denote its derivatives. The printed conjecture, attributed to B. Sh, is
\[
 \#_r\!\left[(k-1)(p')^2-kpp''\right]+\#_r p>0.
\]
The record does not define or quantify $k$. In particular it does not say whether $k=\deg p$, a fixed parameter, or an independently quantified variable. No such relation is silently imposed.

The notation $\#_r q$ evidently concerns real roots of a polynomial $q$, but the supplied record does not define its counting convention. For nonzero polynomials, either the distinct-root count or the count with multiplicity is nonnegative, and positivity of their sum means that at least one of
\[
 p(x),\qquad (k-1)(p'(x))^2-kp(x)p''(x)
\]
has a real zero. This explains the strict inequality without selecting a multiplicity convention. The second expression can be identically zero for some choices; the root-count convention for that case is not supplied. The source also does not exclude constant polynomials, which have even degree zero, or otherwise state a minimum degree.

Accordingly, the exact formula is retained, but a fully quantified numerical conjecture requires recovering $k$, the degree range, and the zero-polynomial root-count convention from the original context.
\subsection{Short English statement}
The conjecture asserts that an even-degree real polynomial or the displayed derivative combination must have a real root. The supplied assessment labels the conjecture false: it attributes counterexamples in every even degree at least four to Katkova--Tyaglov--Vishnyakova and a 2025 classification to Ma--Ma, with nine disjoint cases where the inequality holds and four where it fails. The record supplies neither the cases nor complete references, and leaves the parameter $k$ undefined, so those resolution claims are retained as attributed rather than asserted as a verified theorem for a guessed formulation.
\subsection{Sources}
\begin{itemize}
\item UnsolvedMath permanent ID 2200019, source code catalogue entry, ``Problems Around Polynomials --- Conjecture 12.'' Registry curation and metadata: Ulam AI and UnsolvedMath Contributors, CC BY 4.0; updated 2026-08-06.
\item Shapiro, \emph{Problems Around Polynomials: The Good, The Bad and The Ugly} (2015), source-order item 19, Conjecture 12 as printed, attribution B. Sh: \url{https://arxiv.org/abs/1503.05295}. Supplied source-TeX extraction accessed 2026-07-29. Original metadata presents an open item; status/release and rights review are NEEDS\_REVIEW.
\item Katkova--Tyaglov--Vishnyakova and Ma--Ma (2025), as named in the supplied resolving assessment. No titles, venues, locators, counterexample formulas, or classification cases are supplied; none are invented.
\end{itemize}
\end{document}
